\documentclass[A4]{article}
\usepackage{amsmath, amssymb, braket}
% 美文書入門 p.74 5.2数式用のフォント
\usepackage{lmodern}
\usepackage{bm}
\usepackage{quantikz}
\usepackage{hyperref}
\hypersetup{
    colorlinks=true,
    linkcolor=blue,
    urlcolor=cyan,
    pdftitle={Quantum Circuit for First Quantization Hamiltonian Simulation Handling Coulomb Potential Pole},
    pdfauthor={Hideo Takahashi},
    pdfsubject={Research Paper},
    pdfkeywords={Quantum Computer, First Quantization, ab initio, Simulator, Coulomb Potential, Quantum Chemistry},
    bookmarksnumbered=true,
    bookmarksopen=true
}

\usepackage[backend=biber, style=numeric, maxnames=5, doi=false, url=true, eprint=true]{biblatex}
\addbibresource{quantum.bib}
\makeatletter
\let\mkbibemph=\mkbibitalic
\let\blx@imc@mkbibemph=\blx@imc@mkbibitalic
\makeatother
\setlength{\biblabelsep}{0.5em}

\usepackage{paper2025}
\begin{document}
\title{Development and error analysis of
a quantum circuit for first quantization Hamiltonian simulation 
of a one and two dimensional hydrogen atom model }
\author{Hideo Takahashi}
\date{\today}
\maketitle
\begin{abstract}
We show a quantum circuit of a chemical simulator for Hamiltonian simulation
based on first quantization which can run on quantum computers with limited
number of qubits. Handling of the Coulomb potential pole is designed and tested.
We also show a method to choose simulation parameters to optimize simulation
accuracy within a given number of qubits for the circuit. Using the uniformly
controlled rotation (UCR) circuit, which can be seen as a variation of quantum
read only memories (QROMs), the Coulomb energy calculation can be implemented
with fewer number of qubits and gates compared to a circuit based on arithmetic
gates when the model qubit count is small, such as in a two-dimensional (2-D)
Hydrogen model.  This enabled first quantization Hamiltonian simulation of 2-D
Hydrogen using only 26 qubits.  Regarding the pole of the Coulomb potential, we
show that replacing the potential value at the pole by a finite value
H(0)=-4/δq, where δq is the spatial discretization grid size, will lead to good
simulation results. When the qubit count for the overall circuit is given, there
is a trade-off between the simulated box edge length L and discretization
resolution δq-1. We show a method to minimize simulation error under this
trade-off for a given wave function.  Using values determined by this method,
good simulation results with low level of artifacts and aliases were obtained
with a coarse 6-qubit spatial coordinate grid.
\end{abstract}

\section{Introduction}
Quantum computing is expected to be a powerful tool for simulating quantum
systems \cite{feynman1982simulating,AbramsAndLloyd1997.PhysRevLett.79.2586}. In particular, quantum chemistry
simulations using quantum computers have been actively studied in recent years

\cite{Aspuru-Guzik2005.10.1126/science.1113479, kassal2008.pnas.0808245105, Whitfield2011, Reiher2017, Babbush2018,
Motta2020, Cao2019, McArdle2020, Kandala2017, OMalley2016, Kandala2019,
Kivlichan2018, Grimsley2019, Sugisaki2020, Takahashi2023}. Most of these
studies are based on the second quantization formalism, where the electronic
structure is represented using a finite set of basis functions, and the
Hamiltonian is expressed in terms of creation and annihilation operators for
these basis functions. However, the second quantization approach requires a
number of qubits proportional to the number of basis functions, which can be
large for accurate simulations. On the other hand, the first quantization
formalism represents the wave function directly in real space, allowing for
more compact representations for certain systems. First quantization Hamiltonian
simulation has been studied in several works \cite{Zalka1998, Kassal2006,
Babbush2018first, Kivlichan2017, Berry2019first, Childs2021first}. However,
implementing the Coulomb potential term in the Hamiltonian can be challenging due to
the singularity at the origin. In this paper, we present a quantum circuit for
first quantization Hamiltonian simulation that efficiently handles the Coulomb
potential pole and optimizes simulation parameters for a given qubit budget. We
demonstrate our approach using a two-dimensional Hydrogen model, achieving
accurate results with a limited number of qubits.
\section{Quantum Circuit for First Quantization Hamiltonian Simulation}
In first quantization, the Hamiltonian for a single electron in a Coulomb
potential is given by
\begin{equation}
H = -\frac{1}{2} \nabla^2 - \frac{1}{r},
\end{equation}
where the first term represents the kinetic energy and the second term
represents the Coulomb potential. To simulate the time evolution operator
$U(t) = e^{-iHt}$, we use the Trotter-Suzuki decomposition to separate the kinetic and
potential terms:
\begin{equation}
U(t) \approx \left( e^{-iT\Delta t} e^{-iV\Delta t} \right)^n,
\end{equation}
where $T = -\frac{1}{2} \nabla^2$, $V = -\frac{1}{r}$, $\Delta t = t/n$, and $n$ is the number of Trotter steps.
\subsection{Kinetic Energy Term}
The kinetic energy term can be implemented using the quantum Fourier transform
(QFT) to switch to momentum space, apply the kinetic energy operator, and then
transform back to position space. The circuit for the kinetic energy term is
shown in Figure~\ref{fig:kinetic}.
\begin{figure}[h]
\centering
graphics come here.
\caption{Quantum circuit for the kinetic energy term using QFT.}
\label{fig:kinetic}
\end{figure}
\subsection{Coulomb Potential Term}
The Coulomb potential term $V = -\frac{1}{r}$ has a singularity at $r=0$. To handle this,
we approximate the potential near the singularity by assigning a finite value
$V(0) = -\frac{4}{\delta q}$, where $\delta q$ is the spatial discretization grid size. This value is
determined by calculating the weighted average of $V(r)$ near the pole, assuming a specific wave function.
The circuit for the potential term is implemented using a uniformly controlled
rotation (UCR) circuit, which allows for efficient implementation of the potential
term with fewer qubits and gates. The circuit is shown in Figure~\ref{fig:potential}.
\begin{figure}[h]
\centering
graphics come here.
\caption{Quantum circuit for the Coulomb potential term using UCR.}
\label{fig:potential}
\end{figure}
\section{Optimization of Simulation Parameters}
When the total number of qubits available for the simulation is limited, there
is a trade-off between the simulated box edge length $L$ and the discretization
resolution $\delta q^{-1}$. To minimize the simulation error under this trade-off, we propose a method to choose optimal values for $L$ and $\delta q$ based on the
given wave function. The optimization process involves evaluating the simulation
error for different combinations of $L$ and $\delta q$, and selecting the pair
that minimizes the error while satisfying the qubit budget constraint.
\section{Results}
We tested our quantum circuit and optimization method using a two-dimensional
Hydrogen model. The simulation was performed using a total of 26 qubits, with
6 qubits allocated for the spatial coordinate. The results showed good agreement
with the expected energy levels and wave functions, with minimal artifacts and
aliases. Figure~\ref{fig:results} shows the comparison between the simulated and
exact wave functions.
\begin{figure}[h]
\centering
graphics come here.
\caption{Comparison of simulated and exact wave functions for the 2-D Hydrogen model.}
\label{fig:results}
\end{figure}
\section{Conclusion}
We have presented a quantum circuit for first quantization Hamiltonian
simulation that efficiently handles the Coulomb potential pole using a UCR
circuit. We also proposed a method to optimize simulation parameters for a given
qubit budget. Our approach was demonstrated using a two-dimensional Hydrogen
model, achieving accurate results with a limited number of qubits. This work
paves the way for more efficient quantum simulations of quantum systems using
first quantization methods.

\printbibliography[heading=bibintoc, title={参考文献}]
\end{document}